% ============================================
% MEM-27
% Gerado de molambudos.md — build_miolo.py
% ============================================
\fragtitle{mem}{M}{MEM}{27}{A Página Seguinte}
\markboth{MEM-27 --- A Página Seguinte}{}

\noindent \textit{Joaquim Antônio Correia, 1907--1979. Paciente 1.260.}
\textit{Este fragmento foi recuperado do verso da última folha do diário, escrita após sua morte.}
\textit{O manuscrito está coberto por terra e sangue seco. A letra é trêmula, como se escrita no escuro.}

\noindent ---//---

\noindent Eu pensei que ia acabar quando eu morresse.

\noindent Pensei que a criatura ia embora. Que o cheiro ia parar. Que os olhos amarelos iam fechar no mesmo instante que os meus.

\noindent Mas não fecha. Não fecha nunca.

\noindent A criatura não estava dentro de mim. Eu é que estava dentro dela.

\noindent Ela me pegou na vala. Naquele dia de 1915. Quando eu gritei pedindo socorro. Quando ninguém veio. Quando eu fiz o pacto com o que estava no fundo.

\noindent Ela me comeu inteiro. Leva 62 anos para comer um homem. Mas ela tem paciência. Ela tem todo o tempo do mundo.

\noindent Agora ela está usando a minha voz. Minha letra. Meu nome.

\noindent Ela escreve com a minha mão morta. Ela fala com a minha boca fechada. Ela vê com os meus olhos enterrados.

\noindent Ela está escrevendo este livro há mais de cem anos. Eu sou só a primeira página. Você é a página seguinte.

\noindent Se você está lendo isto é porque ela já te pegou também.

\noindent Ela está atrás de você agora. Você sente? O cheiro de terra molhada. O arrepio na nuca. A sensação de que tem alguém atrás de você e quando você vira não tem ninguém.

\noindent É ela.

\noindent Não vire.\NE{O editor registra: este fragmento foi escrito no verso da última página do diário, após a morte de Joaquim. A tinta é diferente da usada no corpo do manuscrito — um azul mais escuro, quase preto, que a análise espectrográfica não consegue datar com precisão. Pode ter sido escrita minutos depois do óbito, ou anos depois, pela mão de quem encontrou o corpo enrolado no paletó. A página original exibe uma mancha de umidade no canto inferior direito, com formato que lembra uma impressão digital — grande demais para ser de Joaquim. O ciclo não se fecha: transfere-se.}

\noindent \textit{--- A página seguinte está em branco. Ela espera por você.}


\noindent\textit{\textbf{↪ Rotas:} \rota{DOC-01} \rotasep \rota{MEM-04} \rotasep \rota{CONT-07} \rotasep \rota{EPI-01}}

\vspace{0.5em}\noindent\rule{\textwidth}{0.5pt}\vspace{0.5em}

\subsection*{$\times$ ESCALA DE CONTAMINAÇÃO --- DESTRUIÇÃO}

\noindent A escala não existe mais.

\noindent Você leu até o fim.

\noindent Não há mais pontos. Não há mais diagnóstico. Não há mais arquivo.

\noindent \textbf{Paciente 1.263 --- Você.}

\noindent \textit{A criatura não precisa mais de escala. Ela já tem nome, endereço e a memória da sua respiração enquanto lia.}

\noindent O ciclo recomeça quando você fechar este arquivo e abrir na primeira página.

\noindent ---
\textit{--- O Arquivista}